\section{Polynomial representation of raw singleton blocks}
\label{sec:persistent-elimination}

The existing word arena reconstructs substituted concatenations after every
acyclic elimination batch. Its one-batch polynomial bound does not, by
itself, give a polynomial bound for an arbitrary number of batches. We now
give a different representation with a bound for an entire sequence of raw
singleton eliminations. The representation has a tested reference
implementation, but it is not integrated into the knot-group producer or
the independent source-replay checker.

SLPs and their use in group presentations are established: for example,
Theorem 4.10 of Ennes and Maria's \emph{Compressed data structures for
Heegaard splittings} (arXiv:2507.11406v2) converts compressed relators to
a short ordinary presentation by introducing SLP nonterminals as Tietze
generators~\cite{ennesmaria2025}. The contribution here is an explicit representation, size
recurrence and prototype for the cumulative cost of a complete raw
singleton-elimination block. No claim of literature-wide priority is
needed for this refinement.

\subsection{Signed circuits and mutable generator bindings}

Fix a finite generator alphabet $X$ of size $r_0$. A signed handle $v$
references a circuit vertex; $-v$ represents the formal inverse word of
$v$. Vertex zero represents the empty word. There is one terminal vertex
for each positive generator, and immutable concatenation vertices
$c(u,v)$. Consequently
\[
 \operatorname{val}(-c(u,v))
 =\operatorname{val}(-v)\operatorname{val}(-u).
\]
Inversion requires no additional circuit vertices. Initially the circuit
is a binary SLP. Later, a terminal for an eliminated generator receives
one binding to a signed handle. An unbound terminal denotes its live
generator, while a bound terminal denotes the current value of its target.
All dependencies, including bindings, must form a directed acyclic graph.
Concatenation vertices retain their original children for the whole block;
only previously unbound generator terminals can acquire bindings.

The represented words are raw words. In particular, the adjacent letters
$x x^{-1}$ contribute two occurrences of $x$, even though their product
is trivial in the free group. This restriction makes a zero occurrence
count equivalent to the absence of a path to the corresponding live
terminal. It is essential to the proof below.

The current presentation retains the original list of relator slots.
Eliminating a generator clears exactly its chosen donor slot. Every other
root remains a reference to its existing circuit, so a new generator
binding applies the substitution to all retained relators simultaneously.
Historical donor vertices can remain allocated. Their continued presence
does not affect the presentation and is included in all size bounds.

\subsection{Position-free extraction of a singleton context}

Let $x$ be live and let a donor root represent
\[
   W=Lx^\varepsilon R,\qquad \varepsilon\in\{1,-1\},
\]
with exactly one occurrence of $x$ or $x^{-1}$. For this one generator,
compute counts in the three-element set $\{0,1,2\}$, where $2$ means
at least two. Counts add with saturation at concatenations, pass through
bindings, and are unchanged by orientation. A topological pass therefore
recognizes the singleton condition exactly. Counts need no expanded-length
integers, substring coordinates or word-equality oracle.

Starting at the donor root, follow its unique path to the live terminal
$x$. At a negative concatenation, reverse the order of the children and
negate their handles. At a binding, follow the bound handle with the
appropriate sign. At each concatenation on the path, exactly one child
contains $x$ and the other contains none. Save this other child as a
signed sibling handle. Siblings before the path form $L$ in descent order;
siblings after the path form $R$ in reverse descent order. This yields an
ordered list of existing signed handles for $RL$.

Combine the list by balanced pairwise concatenation. If there are $q$
nonzero handles, at most $\max(0,q-1)$ new concatenation vertices are
needed, and every path through the new concatenations has length at most
$\lceil\log_2\max(1,q)\rceil$. The balancing is by the number of sibling
handles, not by the expanded lengths of their values. Interning repeated
concatenations is optional and can only reduce allocation.

Install the binding
\begin{equation}\label{eq:persistent-binding}
 x\longmapsto
 \begin{cases}
   (RL)^{-1},&\varepsilon=1,\\
   RL,&\varepsilon=-1.
 \end{cases}
\end{equation}
This is the familiar exact solution of $Lx^\varepsilon R=1$.

\paragraph{Lemma (acyclicity and simultaneous substitution).}
The binding in \eqref{eq:persistent-binding} preserves acyclicity and
applies exactly the corresponding singleton Tietze elimination to all
retained roots.

\begin{proof}
Before the binding is installed, every saved sibling has occurrence count
zero for $x$. Hence its old dependency graph has no path to the live
terminal $x$. The new image consists only of these siblings and new
concatenations. It consequently has no path to $x$. Any newly created
directed cycle would have to use the new edge out of $x$ and then return
from the image to $x$, which is impossible.

Let $\phi$ evaluate the old binding graph as raw words in the live
alphabet, and let $\sigma$ replace $x$ by the raw word specified in
\eqref{eq:persistent-binding}. Evaluation of the enlarged acyclic graph
is $\sigma\circ\phi$: this follows directly by induction over its
dependency order, with the only changed terminal being $x$. In
particular, bindings of earlier eliminated generators update their values
automatically. Deleting the donor equation and eliminating $x$ is an
elementary Tietze transformation, so the presented group is preserved.
\end{proof}

The no-path argument remains valid even when a later image references an
earlier eliminated terminal. If that earlier terminal could lead back to
the current generator, its occurrence count in the proposed image would
be positive and it could not be a saved sibling. No fixed elimination
order on original generator labels is required.

\subsection{The global size theorem}

Write $S_0$ for the initial number of nonzero vertices and $D_0$ for the
largest number of concatenation vertices on a dependency path in the
initial SLP. Suppose that exactly $k\le r_0$ distinct live generators are
eliminated, with no intervening free reductions, new generators, or other
word transformations. Donors may be chosen adaptively, but each must
satisfy the raw singleton condition at the time it is used.

After $i$ eliminations let $D_i$ be the largest number of concatenation
vertices on a path when all generator bindings are ignored, and let
$H_i$ be the longest full dependency path, counting concatenations and
binding edges. At most $i$ bindings lie on any full path: a directed
acyclic path cannot visit a terminal twice. Such a path splits into at
most $i+1$ binding-free segments. Therefore
\begin{equation}\label{eq:persistent-height}
 H_i\le(i+1)D_i+i < (i+1)(D_i+1)=:A_i.
\end{equation}
The singleton descent encounters at most $H_i$ concatenations, so its
sibling count $q_i$ is at most $A_i$. Balanced construction gives the
recurrence
\begin{equation}\label{eq:persistent-recurrence}
 D_{i+1}\le D_i+
       \left\lceil\log_2\max(2,A_i)\right\rceil.
\end{equation}

\paragraph{Theorem (complete raw singleton blocks have polynomial size).}
For every $0\le i\le k$,
\begin{equation}\label{eq:persistent-D}
 D_i\le D_0+6i\log_2(D_0+i+2).
\end{equation}
In particular, putting $T=D_0+k+2$, the number $S_k$ of allocated
nonzero vertices satisfies the explicit bound
\begin{align}\label{eq:persistent-S}
 S_k\le S_*:={}&S_0+\frac{k(k+1)}2(D_0+1)\notag\\
 &\quad+2k(k-1)(k+1)\log_2 T.
\end{align}
The bit length $B_k$ of every represented raw length, including historical
roots and image vertices, satisfies
\begin{equation}\label{eq:persistent-B}
 B_k\le(k+1)\bigl(D_0+6k\log_2 T+1\bigr).
\end{equation}
Thus
\[
 S_k=O\bigl(S_0+k^2(D_0+1)+k^3\log(D_0+k+2)\bigr),
\]
and a whole block has polynomial encoded complexity in the initial
grammar size and rank, even if it uses a linear number of serial phases.

\begin{proof}
Set $T_i=D_0+i+2\ge2$. The claim for $D_0$ is immediate. Suppose that
\eqref{eq:persistent-D} holds at $i$. Since $i+1\le T_i$,
$D_0+1\le T_i$, and $i\le T_i$, we obtain
\[
 A_i\le T_i^2(1+6\log_2 T_i)\le T_i^5.
\]
For the final inequality it suffices to observe that
$\log_2 t\le t-1$ for $t\ge2$, and
$6t-5\le t^3$ on the same range. The latter follows because
$t^3-6t+5$ is positive at $t=2$ and has positive derivative thereafter.
As $T_i^5\ge2$, it follows that
\[
 \left\lceil\log_2\max(2,A_i)\right\rceil
 \le5\log_2 T_i+1\le6\log_2 T_i.
\]
Substitution into \eqref{eq:persistent-recurrence} yields
\[
 D_{i+1}\le D_0+6(i+1)\log_2 T_i
          \le D_0+6(i+1)\log_2 T_{i+1},
\]
completing the induction with the stated constant.

At step $i$, no generator vertex is added and at most $q_i\le A_i$
concatenation vertices are allocated. Sum
\[
 A_i\le(i+1)(D_0+1)+6i(i+1)\log_2 T
\]
over $i=0,\ldots,k-1$. The identities
\[
 \sum_{i=0}^{k-1}(i+1)=\frac{k(k+1)}2,\qquad
 \sum_{i=0}^{k-1}i(i+1)=\frac{k(k-1)(k+1)}3
\]
give \eqref{eq:persistent-S}.

Unfolding a binary/unary acyclic circuit of height $H_i$ gives a tree
with at most $2^{H_i}$ terminal leaves. Inversion does not change length,
and bindings are unary. Hence every represented length has bit length
at most $H_i+1$. Combine this observation with
\eqref{eq:persistent-height} and \eqref{eq:persistent-D} to obtain
\eqref{eq:persistent-B}. Empty values have bit length zero and satisfy
the same bound.
\end{proof}

\subsection{Running time and ordinary SLP export}

For a supplied sequence of donor slots and live generator labels, the
reference algorithm performs a topological pass and a capped occurrence
pass over all historical vertices at each step. The graph has at most
two outgoing edges per vertex. Context extraction and balanced assembly
take $O(q_i)$ further operations. The whole block therefore uses
$O(kS_*)$ graph operations beyond initial input handling. If $M$ is the
original number of relator slots, reading the input and exporting its roots
add $O(S_0+M)$ operations; the total, including final export, is
$O(S_0+M+kS_*)$. Working storage is $O(S_*+M)$, including the fixed
relator-root list and the $k$ binding entries. The exported root list uses
$O(M\log(S_*+r_0+2))$ bits. Input labels may be
renumbered densely once; node identifiers use
$O(\log(S_*+r_0+2))$ bits.

This count requires neither large-integer lengths nor exact string
equality. A deterministic implementation can omit interning and use dense
arrays for all generator and vertex metadata. If interning is desired,
balanced comparison dictionaries give a conservative
$O(kS_*\log^2(S_*+r_0+2))$ bit-operation accounting in the usual
bit-cost random-access model. The Python prototype uses hash dictionaries;
its measured timings do not constitute that deterministic machine bound.

Adaptive donor discovery remains polynomial within this raw-only model:
one may compute the same saturated counts for each live generator and
scan every retained root, adding a factor at most $r_0$ to the graph
passes plus the root-scanning cost. This observation gives no guarantee
that another singleton donor exists, that a chosen order reaches rank
one, or that searching over many alternative orders is cheap.

The representation is not a permanent barrier to existing word kernels.
At an explicit end-of-block boundary, topologically evaluate each vertex
in both orientations. A bound generator becomes an alias to its image;
an unbound generator becomes its positive and negative terminal letters;
a concatenation becomes two ordinary concatenations with opposite child
orders and signs. This exports a standard binary SLP containing at most
$2S_k$ nonzero vertices, with all references in allocation order and no
bindings. The export is linear in circuit size. It is implemented by
\texttt{export\_slp} in the prototype and checked against a separate
literal reconstruction on the finite tests.

A single existing polynomial compressed free-reduction operation
\cite{schleimer2008} can
therefore be charged after the export. The theorem does not cover a
repeated alternation of normalizations and other presentation moves.
Nor does it automatically include torsion-dependent primitive-power
quotients whose donor generator is not a raw singleton. Those operations
require their own growth analysis.

\subsection{A controlled family with exponentially long raw contexts}

Let the generators be $a,b,x_0,\ldots,x_k,z$, and take
\[
 r_0=x_0^{-1}a x_1 b,\qquad
 r_i=x_0^{-1}x_{i+1}\ (1\le i<k),\qquad
 r_k=x_0^{-1}z.
\]
Eliminate $x_i$ from the current value of $r_i$, successively for
$i=0,\ldots,k-1$, and retain $r_k$. After the first elimination,
$x_0=A_1x_1B_1$ with $A_1=a$, $B_1=b$. If after step $j$ it is
$A_jx_jB_j$, the next donor gives
$x_j=A_j^{-1}x_{j+1}B_j^{-1}$, so
\[
 A_{j+1}=A_jA_j^{-1},\qquad B_{j+1}=B_j^{-1}B_j.
\]
These are raw words, so $|A_j|=|B_j|=2^{j-1}$ and the retained relator
has length $2^k+2$. All selected pivots remain singletons. This artificial
presentation family stresses retained contexts and repeated substitutions;
it is not a family of difficult knot diagrams. Free reduction would
remove its deliberately retained cancellation padding.

The benchmark compares the new circuit with the pinned repository's
\texttt{apply\_batch} called with the very same one-entry donor sequence.
Both arms start from the same explicit source relators, clear the same
slots, retain raw words, and disable their expansion method. An identical
eager control measures scheduling variability. There is one excluded
warm-up per size and arm, followed by five randomized A/A/B rounds.
Timing includes source construction and all eliminations; final diagnostic
measurements are excluded. Both arms use caps of two million vertices
and one billion implementation work units, and every run completes.

\begin{center}
\begin{tabular}{r|rrr|rr}
 $k$ & Eager (s) & Control (s) & Persistent (s)
     & Eager vertices & Persistent vertices\\\hline
 8   & 0.000879 & 0.000869 & 0.000363 & 238 & 83\\
 16  & 0.005813 & 0.005988 & 0.001871 & 914 & 279\\
 32  & 0.043468 & 0.045512 & 0.010994 & 3610 & 839\\
 64  & 0.341194 & 0.348591 & 0.067422 & 14378 & 2534\\
 128 & 2.853190 & 3.013364 & 0.395421 & 57418 & 6890
\end{tabular}
\end{center}
The final row gives a measured median speedup of about $7.22$ relative
to the eager arm. Its retained word has length $2^{128}+2$. The new
representation has 128 extra binding entries in addition to its 6,890
vertices, and uses signed handles instead of separate inverse subgraphs;
vertex counts are consequently structural counts, not byte measurements.
Both sides retain historical vertices. Work counters have different
charging conventions and should not be interpreted as portable CPU
instruction counts.

The six test methods include 600 deterministic random presentation traces
checked against independently implemented literal substitution, explicit
negative handles and empty images, repeated-child rejection, forward
bindings, interruption before a binding is published, and the 128-step
capacity example with expansion disabled. The tests check the syntactic
height recurrence and the $2S_k$ export bound as well. These are tests of
word representation and Tietze semantics; no new knot verdict is produced.

\subsection{Consequences and next questions}

The size theorem removes a logarithmic-phase hypothesis from this
restricted raw singleton-elimination component, provided the persistent
representation is used. It does not establish a quasi-polynomial
unknot-recognition algorithm. The implementation work needed before
production use is concrete: add a source-derived replay path; preserve
all presentation-slot and cancellation semantics; implement a compatible
raw donor planner; export exactly at normalization boundaries; and measure
complete discovery and replay on real diagrams. Algebraically, the next
questions are whether the cubic logarithmic size term can be improved,
whether suitable normalization steps preserve a comparable amortized
bound, and which topologically meaningful diagram classes admit enough
raw singleton eliminations to reach a checked terminal presentation.
